\section{Arbitrary convex bodies and consequences}
\label{sec:approximation}

The passage from finite strip bodies to arbitrary convex bodies uses an
approximation in the polar. This keeps the momentum factor inside the
desired polar and enlarges the position factor by an arbitrarily small
amount.

\begin{lemma}[Finite polar approximation]
\label{lem:polar-approximation}
Let $K\subset\R^n$ be an origin-symmetric convex body. For every
$\alpha>0$, there are finitely many vectors
$b_1,\ldots,b_m\in\R^n$ spanning $\R^n$ such that
\[
 K'=\{x\in\R^n:|\ip{b_j}{x}|\le1,\quad 1\le j\le m\}
\]
satisfies
\begin{equation}\label{eq:polar-approximation}
 K\subset K'\subset(1+\alpha)K,
 \qquad (K')^\circ\subset K^\circ.
\end{equation}
\end{lemma}

\begin{proof}
First, $0\in\interior K$. Indeed, if the Euclidean ball $B(y,r)$
is contained in $K$, then so is $B(-y,r)$, and their midpoints contain
$B(0,r)$. Since $K$ is also bounded, there are $r,R>0$ with
$B(0,r)\subset K\subset B(0,R)$. The definition of the polar gives
\[
 B(0,R^{-1})\subset K^\circ\subset\overline{B(0,r^{-1})}.
\]
Thus $K^\circ$ is a compact, origin-symmetric convex body with nonempty
interior.

We will use the bipolar identity $K^{\circ\circ}=K$. To recall its
justification here, if $x\notin K$, separation gives a nonzero vector
$v$ such that
\[
 \ip{v}{x}>h_K(v),
 \qquad h_K(v)=\max_{y\in K}\ip{v}{y}>0.
\]
The normalized vector $v/h_K(v)$ belongs to $K^\circ$ and pairs with
$x$ to a value greater than $1$, so $x\notin K^{\circ\circ}$.
The reverse inclusion follows directly from the definition. The same
argument applies to every convex body containing the origin in its
interior.

Put $\lambda=1+\alpha$ and
\[
 \rho=\min_{|v|=1}h_{K^\circ}(v)>0.
\]
Choose $d>0$ with $d\le\alpha\rho/(1+\alpha)$, and choose a finite
$d$-net $\{b_1,\ldots,b_m\}$ in the compact set $K^\circ$, with all
net points in $K^\circ$. Let
\[
 P=\operatorname{conv}\{\pm b_1,\ldots,\pm b_m\}.
\]
Symmetry of $K^\circ$ gives $P\subset K^\circ$. For every unit vector
$v$, a maximizer of $\ip{v}{p}$ over $p\in K^\circ$ lies within
distance $d$ of a net point. Consequently
\[
 h_P(v)\ge h_{K^\circ}(v)-d
       \ge\lambda^{-1}h_{K^\circ}(v).
\]
The characterization of inclusion by support functions, which follows
from convex separation, now gives
\[
 \lambda^{-1}K^\circ\subset P\subset K^\circ.
\]
In particular, $P$ has interior, so its generating vectors span
$\R^n$. Taking polars and using the bipolar identity yields
\[
 K\subset P^\circ\subset\lambda K,
 \qquad (P^\circ)^\circ=P\subset K^\circ.
\]
Finally, the definition of $P$ shows that $P^\circ$ is exactly the
finite strip body $K'$ in the statement.
\end{proof}

\begin{proof}[Proof of Theorem~\ref{thm:main}]
Fix $0<c<4$. Choose $\alpha>0$ such that $\lambda c<4$, where
$\lambda=1+\alpha$, and choose one body $K'$ from
Lemma~\ref{lem:polar-approximation}. Proposition~\ref{prop:finite-strips}
gives a smooth symplectic embedding
\[
 e:B^{2n}(\lambda c)\longrightarrow
 \interior K'\times\interior\bigl((K')^\circ\bigr).
\]
If one body is contained in another, every interior ball in the first
is contained in the second. Thus Equation~\eqref{eq:polar-approximation}
implies
\[
 \interior K'\subset\lambda\interior K,
 \qquad \interior\bigl((K')^\circ\bigr)\subset\interior K^\circ,
\]
even when their boundaries meet. Regard $e$ as an embedding into
$\lambda\interior K\times\interior K^\circ$. Define
\[
 D_{\sqrt\lambda}(z)=\sqrt\lambda\,z,
 \qquad Q_\lambda(q,p)=(q/\lambda,p).
\]
The composition
\begin{equation}\label{eq:general-rescaling}
 Q_\lambda\circ e\circ D_{\sqrt\lambda}:
 B^{2n}(c)\longrightarrow U_K
\end{equation}
is a smooth embedding. Since
\[
 D_{\sqrt\lambda}^*\omega_0=\lambda\omega_0,
 \qquad Q_\lambda^*\omega_0=\lambda^{-1}\omega_0,
\]
its pullback of $\omega_0$ equals $\omega_0$. This proves
$c_G(U_K)\ge4$.

For the upper bound, we use the supporting-slab construction of
Artstein-Avidan, Karasev, and Ostrover
\cite[Remark~4.2]{AAKO2014}. Choose $q_0\in K$ and $p_0\in K^\circ$
with $\ip{q_0}{p_0}=1$. Such a pair is obtained by maximizing any
nonzero linear functional on $K$ and normalizing its maximum to $1$.
Complete $q_0$ to a basis of $\R^n$ by vectors in $\ker p_0$, and let
$A$ be the matrix of this basis. The linear change
\[
 Q=A^{-1}q,\qquad P=A^Tp
\]
is symplectic and has first coordinates
$Q_1=\ip{p_0}{q}$ and $P_1=\ip{q_0}{p}$.
Both functionals are nonzero, so symmetry, polarity, and the interior
condition give $|Q_1|<1$ and $|P_1|<1$ on $U_K$.
Thus this change embeds $U_K$ into
$(-1,1)^2\times\R^{2n-2}$, with the square in its first conjugate
coordinate pair.

For completeness, this open square is area-preservingly diffeomorphic
to the disk of area $4$. Put $R=2/\sqrt\pi$,
$h(X)=\sqrt{R^2-X^2}$ for $-R<X<R$, and
\[
 H(X)=-1+\int_{-R}^{X}h(s)\,ds.
\]
Since $\int_{-R}^{R}h(s)\,ds=2$ and $H'=h>0$,
$H$ is a smooth diffeomorphism from $(-R,R)$ onto $(-1,1)$.
The map
\[
 (u,v)\longmapsto
 \bigl(X=H^{-1}(u),\ Y=v\,h(H^{-1}(u))\bigr)
\]
has image $B^2(4)$ and Jacobian $X'(u)h(X)=1$.
Applying it to the first conjugate pair gives a symplectic embedding
$U_K\hookrightarrow B^2(4)\times\R^{2n-2}$.
Gromov's nonsqueezing theorem \cite{Gromov1985} therefore bounds the
capacity of every ball in $U_K$ by $4$. Combined with the lower bound,
this proves $c_G(U_K)=4$.
\end{proof}

\begin{remark}[Order of choices]
\label{rem:parameter-order}
For each prescribed $0<c<4$, first choose $\alpha>0$ so that
$\lambda c<4$, where $\lambda=1+\alpha$. Then fix the finite body
$K'$ and its $m$ defining constraints. Choose $\eta>0$ so that
$(1+\eta)\lambda c<4$, then choose $k$ large enough that
\[
 \frac{C_{K'}}k+m\epsilon_k<\frac{\eta\pi}{4}.
\]
This is possible because the constraint list is fixed before $k$
varies. Proposition~\ref{prop:finite-strips}, applied at capacity
$\lambda c$, and the final rescaling then give the required ball.
Each prescribed capacity is obtained by one finite construction; no
limit of embeddings is required.
\end{remark}

\begin{corollary}[Symmetric Mahler inequality]
\label{cor:mahler}
For every integer $n\ge1$ and every origin-symmetric convex body
$K\subset\R^n$,
\[
 \vol_n(K)\,\vol_n(K^\circ)\ge\frac{4^n}{n!}.
\]
\end{corollary}

\begin{proof}
For $n=1$, write $K=[-a,a]$ with $a>0$. Then
$K^\circ=[-a^{-1},a^{-1}]$, and the product of their lengths is $4$.
Assume henceforth that $n\ge2$.

We first justify that taking interiors does not change the volumes in
the statement. If a convex body $C\subset\R^n$ contains $B(0,r)$,
then convexity gives, for $0<t<1$,
\[
 tC+B(0,(1-t)r)\subset C.
\]
Hence $tC\subset\interior C$, and
\[
 t^n\vol_n(C)\le\vol_n(\interior C)\le\vol_n(C).
\]
Letting $t\uparrow1$ proves
$\vol_n(\interior C)=\vol_n(C)$. Apply this to $K$ and $K^\circ$,
both of which contain a neighborhood of the origin. The product formula
for Lebesgue measure then gives
\begin{equation}\label{eq:polar-product-volume}
 \vol_{2n}(U_K)=\vol_n(K)\,\vol_n(K^\circ).
\end{equation}

For every $0<c<4$, Theorem~\ref{thm:main} supplies a smooth symplectic
embedding $f:B^{2n}(c)\to U_K$. Taking the $n$th exterior power of
$f^*\omega_0=\omega_0$ shows that $\det Df=1$, so the change of
variables formula gives
\[
 \frac{c^n}{n!}
 =\vol_{2n}\bigl(B^{2n}(c)\bigr)
 =\vol_{2n}\bigl(f(B^{2n}(c))\bigr)
 \le\vol_{2n}(U_K).
\]
Here the first equality follows from the Euclidean ball volume
$\pi^nR^{2n}/n!$ at radius $R=\sqrt{c/\pi}$. Combining this with
Equation~\eqref{eq:polar-product-volume} and letting $c\uparrow4$ proves
the claimed inequality.
\end{proof}

\begin{remark}[Sharpness]
The cube $[-1,1]^n$ has volume $2^n$. Its polar is
$\{p:\sum_{j=1}^n|p_j|\le1\}$, whose intersection with each coordinate
orthant is a simplex of volume $1/n!$. These $2^n$ simplices meet only
along faces, so the polar has volume $2^n/n!$. Their volume product is
therefore $4^n/n!$, showing that the constant in
Corollary~\ref{cor:mahler} is sharp.
\end{remark}

\begin{corollary}[A sufficient packing criterion]
\label{cor:polar-packings}
Let $n\ge3$ and $k\ge1$ be integers, and let $K\subset\R^n$ be an
origin-symmetric convex body. Suppose that $R_1,\ldots,R_k>0$ satisfy
\[
 \sum_{i=1}^k R_i^n<4^n,
 \qquad R_i+R_j<4\quad(1\le i<j\le k).
\]
Then there are symplectic embeddings of open neighborhoods of the closed
balls $\overline{B^{2n}(R_1)},\ldots,\overline{B^{2n}(R_k)}$ into $U_K$
with pairwise disjoint images. The pairwise condition is vacuous for $k=1$.
\end{corollary}

\begin{proof}
Strictness and finiteness permit a choice of $0<c<4$ such that
$\sum_i R_i^n<c^n$ and $R_i+R_j<c$ for every $i<j$.
The ball-packing theorem \cite[Theorem~1.1]{OpenAIBallPackings2026}
embeds these closed balls into $B^{2n}(c)$, with embeddings defined on
neighborhoods. Their compact images are pairwise disjoint and lie in the
open target, so the neighborhoods can be shrunk to have pairwise disjoint
images contained in $B^{2n}(c)$. Composing with the embedding in
Theorem~\ref{thm:main} proves the assertion.
\end{proof}
